% \section{Related Work}
% \label{sec:related}
% \paragraph{LLM unlearning and its limitations.}
% Machine unlearning removes specific knowledge or behaviors from trained models via gradient ascent~\citep{jang2023knowledge, yao2024large}, influence functions~\citep{koh2017understanding}, or preference optimization~\citep{rafailov2023direct}, since exact retraining~\citep{bourtoule2021machine} is prohibitive for LLMs.
% Benchmarks such as TOFU~\citep{maini2024tofu} and WMDP~\citep{li2024wmdp} evaluate these methods on standalone models under single-turn, fixed-prompt settings, though concurrent work shows these evaluations can be unreliable under stochastic decoding~\citep{lynch2024eight, shi2024muse}.
% Two recent works extend unlearning to agents:~\cite{wang2026agentic} jointly unlearn entities from parameters and vector-database memory via Synchronized Backflow Unlearning, while Sanyal and Mandal~\citep{sanyal2025agents} achieve inference-time unlearning through multi-agent routing without weight updates.
% Our work differs in that the unlearning target is \emph{behavioral} (toxic influence propagation through emergent conversation state) rather than factual (entity knowledge in explicit memory), and our evaluation uses counterfactual multi-turn rollouts capturing relapse under re-exposure---a failure mode that single-turn probes cannot detect.

% \paragraph{Toxicity propagation and agent security.}
% The Toxicity Echo Effect~\citep{takacstoxicity} shows that LLM agents mirror interlocutor toxicity in dyadic settings; we extend this to multi-agent thread simulations with controlled topology, dose--response analysis, and---crucially---propose interventions rather than only documenting the phenomenon.
% On the security side, ClawWorm~\citep{zhang2026clawworm} demonstrates self-replicating worm attacks across 40{,}000+ agent instances with 100\% payload persistence once written to configuration, following the earlier Morris~II worm~\citep{cohen2024here} in simulated email assistants.
% Their finding that infection is \emph{absorbing}---no self-correction occurs---directly motivates our memory unlearning pathway, which introduces recovery transitions that break this property.
% Indirect prompt injection~\citep{greshake2023not} and adversarial agent manipulation~\citep{gu2024agent, chen2024agentpoison} exploit the flat context trust model where all tokens carry equal authority; our read-side sanitization is related to context privilege isolation~\citep{chen2025struq}, though we target gradual behavioral contamination rather than one-shot instruction hijacking.

% \paragraph{Safety degradation under multi-turn interaction.}
% Safety-aligned models degrade under multi-turn interaction~\citep{chiang2025web, anil2024many}, particularly as tool-using agents.
% Our results offer a mechanistic explanation: degradation arises from \emph{state accumulation}---the growing transcript provides increasing toxic context that overwhelms parametric safety training.
% This motivates our central claim: parameter-level interventions are insufficient when agents operate in a closed loop; backflow-robust safety requires controlling the state, not just the model.

\vspace{-6pt}
\section{Related Work}
\label{sec:related}

\para{Memory and state in LLM agents.} Recent LLM agent architectures increasingly rely on external state: conversation histories, retrieved documents, memory buffers, reflections, and summaries that are reused across turns~\citep{ai2025memorybench,wang2025recursively}. Prior work has shown that such state improves long-horizon interaction, planning, and coordination~\citep{xi2025agentgym}. Generative Agents store natural-language memories, synthesize them into reflections, and retrieve them to guide later behavior~\citep{park2023generative}; Reflexion uses verbal feedback stored in episodic memory to improve subsequent decisions~\citep{shinn2023reflexion}; MemGPT treats long-context interaction as a memory-management problem~\citep{packer2023memgpt}; and AutoGen composes multiple conversable agents through shared interaction patterns~\citep{wu2024autogen}. These systems motivate our setting: memory is not merely a log of past interaction, but an input that shapes future behavior. \textit{However, prior work largely treats memory as a capability mechanism or an object to retrieve from, whereas we study memory as a behavioral safety surface which can appear benign while steering future generations.}

\para{Unsafe context propagation in agentic systems.}
Recent work shows that LLM applications become vulnerable when untrusted context is stored, retrieved, or reused in later model calls. Indirect prompt injection demonstrates that external content can be interpreted as instructions, exploiting the blurred boundary between data and commands in LLM-integrated systems~\citep{greshake2023not}. Memory and retrieval poisoning attacks extend this risk to agentic settings, showing that long-term memory or knowledge bases can be manipulated to steer future agent behavior~\citep{chen2024agentpoison}. Self-propagating agent attacks further show that once malicious payloads enter persistent agent state, they can spread across downstream tools, applications, or agent ecosystems~\citep{cohen2024here, zhang2026clawworm}. Defenses such as structured prompting and context privilege separation aim to prevent untrusted data from being interpreted as privileged instructions~\citep{chen2025struq}. 
\textit{Our work approached this from a complementary perspective, since memory laundering does not require an explicit instruction payload, backdoor trigger, or threshold-crossing toxic artifact.}

\para{From output-level safety and unlearning to state-level mitigation.}
Safety methods commonly intervene on model outputs or model parameters. Interpretability-based techniques detect and block unsafe generations through interventions on model activations~\citep{goyal-etal-2025-breaking,wang2026causaldetox}, while unlearning and preference optimization aim to reduce unsafe knowledge or behaviors inside the model~\citep{jang2023knowledge,rafailov2023direct,yao2024large}. Benchmarks such as TOFU and WMDP evaluate whether specific knowledge or capabilities can be removed from model behavior~\citep{maini2024tofu,li2024wmdp}, and recent work has shown that safety can degrade under long-context, pre-taining, or multi-turn exposure to potentially harmful content~\citep{anil2024many,chiang2025web,xing2025llms}. These approaches are important but insufficient for the failure mode we study. \textit{In memory laundering, it is the external conditioning state that has been transformed by summarization and later reused. This motivates state-level controls over what agents read, what they write, and when memory is updated.}

% \para{LLM unlearning and its limitations.}
% Machine unlearning removes knowledge or behaviors from trained models via gradient ascent~\citep{jang2023knowledge,rafailov2023direct,yao2024large}, influence functions~\citep{koh2017understanding}, or preference optimization~\citep{rafailov2023direct}, avoiding costly exact retraining~\citep{bourtoule2021machine}. Benchmarks such as TOFU~\citep{maini2024tofu,li2024wmdp} and WMDP~\citep{li2024wmdp} evaluate unlearning under standalone, single-turn settings, though recent work shows such evaluations can be unreliable under stochastic decoding~\citep{lynch2024eight, shi2024muse}. Agentic unlearning has begun to address memory-bearing systems: Wang et al.~\citep{wang2026agentic} unlearn entities from parameters and vector-database memory, while Sanyal and Mandal~\citep{sanyal2025agents} use inference-time multi-agent routing. Our target is different: behavioral toxic influence propagated through emergent conversation state, evaluated with counterfactual multi-turn rollouts that expose relapse under re-exposure rather than factual recall in explicit memory.

% \para{Toxicity propagation and agent security.}
% The Toxicity Echo Effect~\citep{takacstoxicity} shows that LLM agents mirror interlocutor toxicity; we extend this phenomenon to controlled multi-agent thread topologies and introduce interventions. Related security work shows that agent state can persist and propagate attacks: ClawWorm~\citep{zhang2026clawworm} demonstrates self-replicating attacks with persistent payloads, following the Morris~II worm~\citep{cohen2024here} in simulated email assistants. Indirect prompt injection~\citep{greshake2023not} and adversarial agent manipulation~\citep{gu2024agent, chen2024agentpoison} exploit flat context trust, while context privilege isolation~\citep{chen2025struq} separates trusted and untrusted inputs. Our focus is complementary: gradual behavioral contamination that survives summarization as classifier-clean memory, rather than one-shot instruction hijacking or explicit payload persistence.

% \para{Safety degradation under multi-turn interaction.}
% Safety-aligned models degrade under multi-turn interaction~\citep{chiang2025web, anil2024many}, especially in agentic or tool-using settings. We provide a state-based mechanism for this degradation: toxic context accumulates through transcript and compressed memory, then re-enters future conditioning contexts and overwhelms purely parametric safeguards. This motivates our central claim that backflow-robust safety requires controlling the agent state, not only the model parameters.

% \noindent \textbf{Positioning.}
% Unlike prior work on filtering, prompt injection, or parameter-level unlearning, we study behavioral influence carried by compressed external agent state: memory can be classifier-clean and payload-free yet still alter downstream behavior. SPG measures this gap.
